\documentclass[11pt]{article}
\usepackage[margin=0.75in]{geometry}
\usepackage{booktabs}
\usepackage{array}
\usepackage{graphicx}
\usepackage{float}
\usepackage{hyperref}
\usepackage{xurl}
\usepackage{xcolor}
\usepackage{caption}
\usepackage{enumitem}
\hypersetup{colorlinks=true, linkcolor=blue, urlcolor=blue}
\setlist[itemize]{topsep=2pt,itemsep=2pt}
\setlist[enumerate]{topsep=2pt,itemsep=2pt}

\title{Synthetic BOAMP Linkage Algorithm Benchmark Report}
\author{Generated from the corrected synthetic benchmark notebook}
\date{2026-08-04}

\begin{document}
\maketitle

\section*{Technical Summary}

Gradient boosting is the strongest current linker on the corrected synthetic benchmark. It reaches
precision \textbf{0.747}, end-to-end recall \textbf{0.235},
and end-to-end pair-F1 \textbf{0.357} across the held-out evaluation grid.

The result is not mainly a classifier-only story. It is a pipeline story: the production candidate
generator exposes only \textbf{10,315} of \textbf{32,050} in-scope truth links
to scoring, a weighted blocking recall of \textbf{32.2\%} and an unweighted
world-average recall of \textbf{0.293}. This is why gradient boosting has
fixed-candidate F1 \textbf{0.738} but end-to-end F1 only
\textbf{0.357}.

Logistic regression recovers essentially the same number of true links as gradient boosting
(7,519 versus 7,517), but accepts far more links
(22,408 versus 10,057). The extra volume lowers
precision from \textbf{0.747} to \textbf{0.336}. The current
weighted composite and Fellegi-Sunter-style baseline remain useful comparators, but they are not the
preferred scorers under this benchmark.

\section*{Gradient Boosting Wins, Mostly By Reducing False Positives}

The four algorithms were evaluated on the same held-out candidate-pair universe. Gradient boosting
accepts fewer links than logistic regression and the composite baseline, while retaining nearly the same
true-positive count as logistic regression. That makes it the best current choice for synthetic
benchmark ranking and the least risky of the tested options for downstream survival-analysis pilots.

\begin{figure}[H]
\centering
\includegraphics[width=0.92\textwidth]{reports/figures/synthetic\_benchmark/v0\_4\_population\_alias\_revision/linkage\_algorithm\_benchmark/pair\_f1\_by\_mode.png}
\caption{Pair-F1 by algorithm and evaluation mode. Fixed-candidate performance is much higher because unreachable truth links are excluded from that denominator.}
\label{fig:pair-f1}
\end{figure}

\begin{table}[H]
\centering
\caption{Aggregate held-out performance by algorithm.}
\label{tab:summary}
\resizebox{\textwidth}{!}{%
\begin{tabular}{lrrrrrr}
\toprule
Algorithm & Predicted links & True positives & Precision & Fixed-candidate recall & End-to-end recall & End-to-end F1 \\
\midrule
Gradient boosting & 10,057 & 7,517 & 0.747 & 0.729 & 0.235 & 0.357 \\
Logistic regression & 22,408 & 7,519 & 0.336 & 0.729 & 0.235 & 0.276 \\
Current weighted composite & 56,017 & 7,189 & 0.128 & 0.697 & 0.224 & 0.163 \\
Fellegi-Sunter style & 42,040 & 4,739 & 0.113 & 0.459 & 0.148 & 0.128 \\
\bottomrule
\end{tabular}%
}
\end{table}

Compared with the current weighted composite, gradient boosting improves precision by
\textbf{5.8x} and end-to-end F1 by
\textbf{2.2x}.

\section*{Candidate Generation Is The Binding Constraint}

End-to-end performance is much lower than fixed-candidate performance because most truth links never
enter the candidate set. The evaluation run scored \textbf{876,190} candidate pairs over
\textbf{32,050} in-scope truth links, but only \textbf{10,315} truth links
were reachable by the production candidate generator.

\begin{table}[H]
\centering
\caption{Candidate-generation reachability by evaluation scenario.}
\label{tab:blocking}
\resizebox{\textwidth}{!}{%
\begin{tabular}{lrrrrrrr}
\toprule
Scenario & Worlds & Candidate pairs & Truth links & Reachable truth & Mean blocking recall & Min & Max \\
\midrule
central\_provisional & 4 & 81,457 & 6,435 & 2,183 & 0.340 & 0.307 & 0.359 \\
difficult & 4 & 175,881 & 3,723 & 851 & 0.231 & 0.196 & 0.267 \\
easier & 4 & 142,948 & 11,526 & 4,449 & 0.387 & 0.350 & 0.402 \\
moderate & 4 & 81,457 & 6,435 & 2,183 & 0.340 & 0.307 & 0.359 \\
stress & 4 & 394,447 & 3,931 & 649 & 0.166 & 0.157 & 0.182 \\
\bottomrule
\end{tabular}%
}
\end{table}

The stress scenario makes this clearest: it creates 386,410 candidate pairs but exposes only 648 of
4,030 truth links to scoring. No scoring model can recover links that were never generated. Therefore
the next major gain should come from candidate generation, not only from another classifier.

\section*{Scenario Results Show Robust Ranking But Different Difficulty}

Gradient boosting ranks first in every evaluated scenario. The ranking is stable, but absolute
performance changes sharply as the benchmark gets harder.

\begin{figure}[H]
\centering
\includegraphics[width=0.92\textwidth]{reports/figures/synthetic\_benchmark/v0\_4\_population\_alias\_revision/linkage\_algorithm\_benchmark/scenario\_pair\_f1.png}
\caption{Held-out end-to-end pair-F1 by scenario. Gradient boosting leads throughout, but stress and difficult scenarios lower all algorithms.}
\label{fig:scenario-f1}
\end{figure}

\begin{table}[H]
\centering
\caption{Mean held-out performance by scenario and algorithm.}
\label{tab:scenario}
\resizebox{\textwidth}{!}{%
\begin{tabular}{lrrrrr}
\toprule
Scenario & Algorithm & Mean F1 & Mean precision & Mean recall & SD F1 \\
\midrule
central\_provisional & Gradient boosting & 0.387 & 0.778 & 0.258 & 0.019 \\
central\_provisional & Logistic regression & 0.354 & 0.561 & 0.259 & 0.020 \\
central\_provisional & Current weighted composite & 0.223 & 0.204 & 0.247 & 0.020 \\
central\_provisional & Fellegi-Sunter style & 0.172 & 0.183 & 0.163 & 0.018 \\
difficult & Gradient boosting & 0.178 & 0.487 & 0.109 & 0.011 \\
difficult & Logistic regression & 0.123 & 0.127 & 0.120 & 0.017 \\
difficult & Current weighted composite & 0.063 & 0.044 & 0.111 & 0.013 \\
difficult & Fellegi-Sunter style & 0.047 & 0.036 & 0.068 & 0.012 \\
easier & Gradient boosting & 0.462 & 0.881 & 0.314 & 0.019 \\
easier & Logistic regression & 0.426 & 0.711 & 0.304 & 0.017 \\
easier & Current weighted composite & 0.280 & 0.271 & 0.290 & 0.014 \\
easier & Fellegi-Sunter style & 0.215 & 0.242 & 0.193 & 0.015 \\
moderate & Gradient boosting & 0.387 & 0.778 & 0.258 & 0.019 \\
moderate & Logistic regression & 0.354 & 0.561 & 0.259 & 0.020 \\
moderate & Current weighted composite & 0.223 & 0.204 & 0.247 & 0.020 \\
moderate & Fellegi-Sunter style & 0.172 & 0.183 & 0.163 & 0.018 \\
stress & Gradient boosting & 0.083 & 0.226 & 0.051 & 0.022 \\
stress & Logistic regression & 0.043 & 0.032 & 0.066 & 0.008 \\
stress & Current weighted composite & 0.026 & 0.016 & 0.072 & 0.005 \\
stress & Fellegi-Sunter style & 0.019 & 0.012 & 0.043 & 0.002 \\
\bottomrule
\end{tabular}%
}
\end{table}

The easier, central, and moderate settings support useful discrimination among algorithms. The
difficult and stress settings are more diagnostic of failure modes. In stress, gradient boosting still
leads, but end-to-end F1 falls to \textbf{0.122} because blocking recall and candidate ambiguity both
worsen. Central and moderate are identical in this executed run, so they should not be interpreted as
separate difficulty levels until their scenario definitions diverge.

\section*{Mechanical Bias Remains Visible}

The benchmark shows systematic performance differences by identifier quality, candidate count, schema,
and CPV availability. These are expected linkage failure modes, but they matter because survival
analysis could inherit them.

\begin{figure}[H]
\centering
\includegraphics[width=0.92\textwidth]{reports/figures/synthetic\_benchmark/v0\_4\_population\_alias\_revision/linkage\_algorithm\_benchmark/recall\_by\_candidate\_count\_bin.png}
\caption{End-to-end recall by candidate count bin. Crowded candidate lists sharply reduce recall.}
\label{fig:candidate-bin}
\end{figure}

For gradient boosting, end-to-end recall is \textbf{0.552}
when a source has exactly one candidate, but only
\textbf{0.113} when it has sixteen or more candidates.
This is the strongest bias diagnostic: crowded candidate environments are under-linked.

\begin{figure}[H]
\centering
\includegraphics[width=0.92\textwidth]{reports/figures/synthetic\_benchmark/v0\_4\_population\_alias\_revision/linkage\_algorithm\_benchmark/recall\_by\_buyer\_key\_type\_heatmap.png}
\caption{End-to-end recall by buyer key type. Cleaner identifiers remain easier to link.}
\label{fig:buyer-key}
\end{figure}

Identifier quality also matters. Gradient boosting recall is
\textbf{0.278} for RAW\_SIRET truth links versus
\textbf{0.224} for NAME\_FALLBACK. CPV availability
matters as well: recall is \textbf{0.251} when CPV is present
and \textbf{0.088} when CPV is missing.

\section*{Thresholds And Model Setup}

The current composite and Fellegi-Sunter thresholds are frozen. Logistic regression and gradient
boosting thresholds are selected on central calibration worlds 005 and 006 by F1 maximization. Training
uses central worlds 001 through 004; evaluation uses worlds 007 through 010 across central provisional,
easier, moderate, difficult, and stress scenarios.

\begin{table}[H]
\centering
\caption{Algorithm score columns and thresholds used in the executed benchmark.}
\label{tab:thresholds}
\resizebox{\textwidth}{!}{%
\begin{tabular}{lrrr}
\toprule
Algorithm & Score column & Threshold & Threshold source \\
\midrule
Current weighted composite & composite\_score & 0.343167 & production primary threshold \\
Fellegi-Sunter style & fs\_log\_odds & 0 & frozen zero log-odds threshold \\
Logistic regression & logistic\_score & 0.757136 & calibration F1 maximum \\
Gradient boosting & gradient\_boosting\_score & 0.221966 & calibration F1 maximum \\
\bottomrule
\end{tabular}%
}
\end{table}

The supervised models use pair-level scoring features already produced by the project pipeline: text
similarity, CPV similarity, time score, buyer reliability score, gap features, candidate count, schema
family, buyer key type, CPV missingness, and duration missingness indicators. The report does not claim
these features are unbiased; the bias diagnostics above show that several are mechanically important.

\section*{Limitations, Robustness, And Interpretation}

These results are valid for the current synthetic benchmark, not for real BOAMP ground truth. The
synthetic truth labels are generated, so model ranking can still contain generator-mechanism bias. The
report supports the claim that gradient boosting is best among the four tested linkers on this synthetic
benchmark; it does not prove real-world BOAMP linkage accuracy.

The aggregate metrics are corrected to use a global pair key containing scenario, world, and notice
pair. This avoids collisions from repeated synthetic notice IDs across generated worlds. The notebook
was rerun after that correction and all tables in this report use the corrected outputs.

The fixed-candidate view is useful for comparing scorers, but the end-to-end view is the
decision-relevant pipeline metric. If downstream survival analysis needs low false positives, gradient
boosting is the best candidate among the tested models. If it needs high recall, the current candidate
generation must be widened or redesigned.

\section*{Recommended Next Steps}

\begin{enumerate}
\item Use gradient boosting as the primary supervised linker for the next synthetic benchmark round.
\item Keep logistic regression as the interpretable supervised baseline, and keep the current composite
and Fellegi-Sunter style models as frozen reference baselines.
\item Improve candidate generation before treating linkage as production-ready. Test wider temporal
windows, alternate buyer-name blocking, CPV fallback blocking, and high-activity buyer handling.
\item Add a manually labeled real BOAMP validation set before making claims about real BOAMP accuracy.
\item Run survival-analysis sensitivity with at least three linkage layers: current composite, gradient
boosting, and a high-precision conservative subset.
\end{enumerate}

\section*{Further Questions}

\begin{enumerate}
\item How much end-to-end recall is recovered by widening candidate generation without creating too many
false-positive candidates?
\item Are high-activity buyers under-linked enough to bias survival estimates?
\item Can a conservative gradient-boosting threshold provide a high-precision survival layer while a
broader threshold supports recall sensitivity?
\item Do manually labeled real BOAMP pairs preserve the same algorithm ranking observed in synthetic truth?
\end{enumerate}

\section*{Source Artifacts}

\begin{itemize}
\item Benchmark version: \texttt{v0\_4\_population\_alias\_revision}
\item Executed notebook: \path{notebooks/14_linkage_algorithm_benchmark_v0_4.ipynb}
\item Summary table: \path{reports/tables/synthetic_benchmark/v0_4_population_alias_revision/linkage_algorithm_benchmark/summary_metrics.csv}
\item Scenario table: \path{reports/tables/synthetic_benchmark/v0_4_population_alias_revision/linkage_algorithm_benchmark/scenario_summary_metrics.csv}
\item Bias diagnostics: \path{reports/tables/synthetic_benchmark/v0_4_population_alias_revision/linkage_algorithm_benchmark/recall_slices.csv}
\end{itemize}

\end{document}
